\documentclass[12pt,a4paper]{article}

% -------------------------------
% Encoding & Language
% -------------------------------
\usepackage{fontspec}        % Needed for Greek with XeLaTeX
\usepackage{polyglossia}
\setmainlanguage{greek}
\setotherlanguage{english}
\usepackage{graphicx}
\usepackage{subcaption}
\usepackage{float}
\usepackage{booktabs}
\usepackage{siunitx}
\usepackage{tikz}
\usetikzlibrary{positioning, arrows.meta, shapes.geometric}
\usepackage{amsmath}
\usepackage{xcolor}

% Tables
\usepackage{booktabs}
\usepackage{longtable}
\usepackage{siunitx}
\sisetup{
  round-mode          = places,
  round-precision     = 4,
  scientific-notation = false
}


% -------------------------------
% Fonts
% -------------------------------
\setmainfont{Noto Serif}     % Works perfectly with Greek
\setsansfont{Noto Sans}

% -------------------------------
% Page & Graphics
% -------------------------------
\usepackage{graphicx}
\usepackage{float}
\usepackage{geometry}
\geometry{margin=2.5cm}

% -------------------------------
% Title info
% -------------------------------
\title{Milestone Report: HNSW (M1) και DARTH πάνω σε HNSW (M2)}
\author{Δαμιανός Ζούμπος}
\date{Φεβρουάριος 2026}

\begin{document}

% ===============================
% Title Page
% ===============================
\begin{titlepage}
    \begin{center}
        \vspace*{-1cm}
        {\Huge \textbf{Milestone: HNSW + DARTH}}\\[1cm]

        \begin{figure}[H]
          \centering
          \framebox{\parbox{0.55\textwidth}{\centering
            \vspace{0.8cm}
            \includegraphics[width=0.8\linewidth]{unipi.png}
            \vspace{0.8cm}
          }}
        \end{figure}

        

        \normalsize
        \renewcommand{\arraystretch}{1.8}
        \begin{tabular}{c|l|c|l}
            \hline
            \textbf{Αρ.} & \textbf{Ονοματεπώνυμο} & \textbf{Α.Μ.} & \textbf{\selectlanguage{english}E-mail\selectlanguage{greek}} \\
            \hline
            1 & Δαμιανός Ζούμπος      & Ε22056 & \texttt{\selectlanguage{english}d.zoumpos04@gmail.com\selectlanguage{greek}} \\
            \hline
        
        \end{tabular}

        \vfill
        
    \end{center}
\end{titlepage}

% ===============================
% Main Content
% ===============================
% ------------------------------------------------------------
% M1 Report (Overleaf-ready) — HNSW from scratch + baseline evaluation
% ------------------------------------------------------------
\section{Μεθοδολογία Εργασίας (Design $\rightarrow$ Review $\rightarrow$ Υλοποίηση $\rightarrow$ Αποτίμηση)}
\label{sec:method}

Ακολουθείται η εξής διαδικασία, όπως προτείνεται για νέες προσεγγίσεις/επεκτάσεις:
\begin{enumerate}
    \item \textbf{Σχεδίαση:} περιγραφή δομών, ροής αλγορίθμου, παραμέτρων και περιπτώσεων (edge cases).
    \item \textbf{Review:} έλεγχος/συζήτηση της σχεδίασης και πιθανές προσαρμογές.
    \item \textbf{Υλοποίηση:} ανάπτυξη κώδικα βάσει της σχεδίασης.
    \item \textbf{Αποτίμηση:} πειραματικός έλεγχος (Recall@k, QPS/Latency, distance computations) και σύγκριση με βιβλιοθήκες αναφοράς.
\end{enumerate}

Σημείωση: Το παρόν report κρατάει \textbf{αυστηρά διαχωρισμένα} τα σκέλη M1 (HNSW baseline) και M2 (DARTH επέκταση),
ώστε να ενσωματωθούν αργότερα στο κείμενο της πτυχιακής χωρίς σύγχυση.

\subsection{Σχεδίαση HNSW: Δομές και Λειτουργίες}
\label{subsec:hnsw_design}

\subsubsection{Δομές δεδομένων}
Ο δείκτης HNSW ορίζεται από:
\begin{itemize}
  \item \textbf{Layers:} $G_L, G_{L-1}, \dots, G_0$ όπου το $G_0$ περιέχει όλους τους κόμβους.
  \item \textbf{neighbors[v][l]:} λίστα γειτόνων του κόμβου $v$ στο layer $l$.
  \item \textbf{entry point:} κόμβος εκκίνησης της πλοήγησης στο ανώτερο layer.
  \item \textbf{level(v):} μέγιστο layer στο οποίο εμφανίζεται ο κόμβος $v$.
\end{itemize}

\subsubsection{Λειτουργία 1: Ανάθεση επιπέδου (Level Sampling)}
Σε κάθε νέο κόμβο ανατίθεται τυχαίο επίπεδο $level(v)$ ώστε οι περισσότεροι κόμβοι να ανήκουν στο $G_0$ και λίγοι σε ανώτερα layers.

\subsubsection{Λειτουργία 2: Greedy descent στα ανώτερα layers}
Για κάθε layer $l = L, L-1, \dots, 1$ εκτελείται greedy πλοήγηση (πλάτος 1) για να βρεθεί καλό σημείο εισόδου στο αμέσως χαμηλότερο layer.

Με τον όρο \textbf{καλό σημείο εισόδου} εννοούμε τον κόμβο που βρίσκεται γεωμετρικά πλησιέστερα στο query $q$ μεταξύ των κόμβων του τρέχοντος layer. Συγκεκριμένα: ξεκινώντας από τον τρέχοντα entry point $ep$, επισκεπτόμαστε τους γείτονές του στο layer $l$ και μεταβαίνουμε στον πλησιέστερο στο $q$. Επαναλαμβάνουμε έως ότου καμία μετακίνηση δεν βελτιώνει την απόσταση. Ο τελικός κόμβος χρησιμοποιείται ως αφετηρία για το layer $l-1$. Η greedy descent \emph{δεν} εγγυάται εύρεση του global nearest neighbor στον layer, αλλά δρα ως γρήγορος οδηγός προς τη σωστή περιοχή του χώρου.

\subsubsection{Λειτουργία 3: \texttt{search\_layer} (Best-first/Beam search)}
Στο layer εισαγωγής και στο base layer εκτελείται best-first αναζήτηση με πλάτος $ef$ (στο build: $efConstruction$, στο query: $efSearch$).

\subsubsection{Λειτουργία 4: Επιλογή γειτόνων (Diversification heuristic)}
Από τους υποψήφιους γείτονες επιλέγονται έως $M$ (ή $M_0$ στο layer 0) με heuristic που αποφεύγει υπερβολικά όμοιους/πλεονάζοντες γείτονες.

\subsubsection{Λειτουργία 5: Αμφίδρομες ακμές και pruning}
Οι ακμές είναι αμφίδρομες: κάθε φορά που ο νέος κόμβος $v$ συνδέεται με κόμβο $u$, προστίθεται και η αντίστροφη ακμή $u \to v$. Αυτό μπορεί να προκαλέσει υπέρβαση του ορίου γειτόνων στον $u$.

\textbf{Πώς εφαρμόζεται το pruning:} αν $|\text{neighbors}[u][l]| > M_{\max}$ (όπου $M_{\max} = M_0$ για $l=0$ και $M$ για $l>0$), τότε εκτελείται εκ νέου η διαδικασία neighbor selection πάνω στο διευρυμένο σύνολο γειτόνων, κρατώντας μόνο τους $M_{\max}$ «ποιοτικότερους». Με τη \emph{simple} παραλλαγή κρατάμε τους $M_{\max}$ πλησιέστερους στο $u$. Με την \emph{heuristic} παραλλαγή (του paper), προτιμούνται γείτονες που «ανοίγουν» νέα μονοπάτια στον γράφο — δηλαδή κόμβοι $e$ για τους οποίους $d(u,e) < d(u,r)$ για όλους τους ήδη επιλεγμένους $r$, αποτρέποντας υπερβολικά «πλεονάζοντες'' ακμές.

% ============================================================
% Toy 2D Example (Expanded) — HNSW full pipeline with cases
% ============================================================

\begin{figure}[H]
\centering
\begin{tikzpicture}[
  node/.style={circle, draw, minimum size=0.7cm, font=\small},
  entrynode/.style={circle, draw, minimum size=0.7cm, font=\small, fill=yellow!40},
  >=Stealth
]

% Layer labels
\node[anchor=east, font=\small\itshape] at (-0.5, 4.0) {Layer 2};
\node[anchor=east, font=\small\itshape] at (-0.5, 2.0) {Layer 1};
\node[anchor=east, font=\small\itshape] at (-0.5, 0.0) {Layer 0};

% Horizontal rules
\draw[gray!40, dashed] (-0.2, 1.1) -- (9.5, 1.1);
\draw[gray!40, dashed] (-0.2, 3.1) -- (9.5, 3.1);

% --- Layer 2 (only 0 and 7) ---
\node[entrynode] (L2_0) at (1.0, 4.0) {0};
\node[node]      (L2_7) at (7.5, 4.0) {7};
\draw[<->] (L2_0) -- (L2_7);

% --- Layer 1 ---
\node[node] (L1_0)  at (1.0, 2.0) {0};
\node[node] (L1_3)  at (3.5, 2.0) {3};
\node[node] (L1_5)  at (5.0, 2.0) {5};
\node[node] (L1_7)  at (7.5, 2.0) {7};
\node[node] (L1_11) at (2.2, 2.0) {11};
\draw[<->] (L1_0)  -- (L1_11);
\draw[<->] (L1_0)  -- (L1_3);
\draw[<->] (L1_3)  -- (L1_5);
\draw[<->] (L1_5)  -- (L1_7);
\draw[<->] (L1_3)  -- (L1_7);

% --- Layer 0 (all nodes) ---
\node[node] (L0_0)  at (0.5, 0.0)  {0};
\node[node] (L0_1)  at (1.5, 0.0)  {1};
\node[node] (L0_2)  at (1.0, -1.0) {2};
\node[node] (L0_3)  at (3.5, 0.0)  {3};
\node[node] (L0_4)  at (4.5, 0.5)  {4};
\node[node] (L0_5)  at (4.0, 0.0)  {5};
\node[node] (L0_6)  at (4.2, 0.8)  {6};
\node[node] (L0_7)  at (7.5, 0.0)  {7};
\node[node] (L0_8)  at (7.5, 0.8)  {8};
\node[node] (L0_9)  at (6.5, 0.5)  {9};
\node[node] (L0_10) at (6.2, 0.0)  {10};
\node[node] (L0_11) at (2.5, 0.8)  {11};

\draw (L0_0)--(L0_1); \draw (L0_0)--(L0_2); \draw (L0_0)--(L0_11);
\draw (L0_3)--(L0_5); \draw (L0_3)--(L0_4); \draw (L0_5)--(L0_6);
\draw (L0_5)--(L0_10); \draw (L0_4)--(L0_6); \draw (L0_4)--(L0_9);
\draw (L0_7)--(L0_9); \draw (L0_7)--(L0_10); \draw (L0_8)--(L0_9);
\draw (L0_10)--(L0_9);

% Vertical dotted guides (same node across layers)
\draw[dotted, gray!60] (L2_0) -- (L1_0) -- (L0_0);
\draw[dotted, gray!60] (L2_7) -- (L1_7) -- (L0_7);
\draw[dotted, gray!60] (L1_3) -- (L0_3);
\draw[dotted, gray!60] (L1_5) -- (L0_5);
\draw[dotted, gray!60] (L1_11)-- (L0_11);

\end{tikzpicture}
\caption{Ιεραρχική δομή HNSW για το toy παράδειγμα. Ο κόμβος 0 (κίτρινος) είναι ο αρχικός entry point στο Layer 2. Οι διακεκομμένες κατακόρυφες γραμμές δείχνουν ότι ο ίδιος κόμβος υπάρχει σε πολλαπλά layers. Στο Layer 2 υπάρχουν μόνο οι κόμβοι 0 και 7 (αραιή δομή για γρήγορη πλοήγηση), ενώ το Layer 0 περιέχει όλους τους κόμβους.}
\label{fig:hnsw_graph}
\end{figure}

\subsection{Toy 2D Παράδειγμα (Εμπλουτισμένο): Πλήρης ροή HNSW με όλες τις περιπτώσεις}
\label{subsec:toy2d_expanded}

Σκοπός του toy παραδείγματος είναι να παρουσιαστούν \textbf{όλες οι βασικές περιπτώσεις} της μεθοδολογίας HNSW,
με \textbf{πραγματικούς υπολογισμούς αποστάσεων} και με \textbf{bounded-degree pruning}.
Χρησιμοποιείται L2 squared distance, ώστε οι αριθμητικές πράξεις να είναι απλές και ελέγξιμες.

\subsubsection{Παράμετροι}
Θέτουμε μικρές παραμέτρους για να είναι πλήρως ιχνηλατήσιμη η διαδικασία:
\[
M=2,\quad M_0=4,\quad efConstruction=4,\quad efSearch=6,\quad k=2.
\]
Η λογική είναι:
\begin{itemize}
    \item Στα layers $>0$ κάθε κόμβος έχει έως $M$ γείτονες.
    \item Στο base layer (layer 0) κάθε κόμβος έχει έως $M_0$ γείτονες.
\end{itemize}

\subsubsection{Dataset 2D και προκαθορισμένα levels (ως αποτέλεσμα sampling)}
Για να αναδειχθεί η ιεραρχία, παρουσιάζουμε τα $level(v)$ ως δεδομένα (toy setup):
\begin{center}
\begin{tabular}{c c c}
\toprule
id & σημείο $(x,y)$ & $level(v)$ \\
\midrule
0 & (0.0,0.0) & 2 \\
1 & (1.0,0.2) & 0 \\
2 & (0.2,1.0) & 0 \\
3 & (3.0,0.0) & 1 \\
4 & (3.0,1.0) & 0 \\
5 & (2.6,0.2) & 1 \\
6 & (2.7,0.9) & 0 \\
7 & (5.0,0.2) & 2 \\
8 & (5.0,1.2) & 0 \\
9 & (4.2,0.7) & 0 \\
10 & (4.0,-0.1) & 0 \\
11 & (2.0,2.0) & 1 \\
\bottomrule
\end{tabular}
\end{center}

Άρα:
\[
G_2=\{0,7\},\quad
G_1=\{0,3,5,7,11\},\quad
G_0=\{0,1,2,\dots,11\}.
\]
Αρχικά entry point θεωρείται ο κόμβος 0 (ο πρώτος με μέγιστο επίπεδο).

% ------------------------------------------------------------
% Case 1: Insert node with maxlevel (update entry point)
% ------------------------------------------------------------
\subsubsection{Case 1: Εισαγωγή κόμβου με νέο/μέγιστο επίπεδο (ενημέρωση entry point)}
\label{subsec:case_maxlevel}

Θεωρούμε εισαγωγή του κόμβου 7:
\[
7=(5.0,0.2),\quad level(7)=2.
\]
Επειδή εμφανίζεται στο ανώτερο επίπεδο, αποτελεί σημαντική περίπτωση:
\begin{itemize}
    \item Ο γράφος στο layer 2 έχει ελάχιστους κόμβους, άρα λειτουργεί ως ``γρήγορη πλοήγηση''.
    \item Η εισαγωγή κόμβου με max level μπορεί να οδηγήσει σε \textbf{ενημέρωση του entry point}.
\end{itemize}

\paragraph{Layer 2: σύνδεση 7 με το υπάρχον top node.}
Πριν το 7, στο layer 2 υπάρχει μόνο ο 0. Με $efConstruction=4$ το \texttt{search\_layer} επιστρέφει προφανώς
$W=\{0\}$, άρα συνδέουμε:
\[
7 \leftrightarrow 0 \quad \text{(στο layer 2)}.
\]
Στη συνέχεια, το 7 συνδέεται και στα layer 1,0 (ίδια μεθοδολογία: local search + neighbor selection).
Στο report, κρατάμε ως βασικό μήνυμα ότι:
\begin{quote}
\textit{Με την εισαγωγή κόμβου υψηλού επιπέδου, ο entry point μπορεί να ενημερώνεται ώστε να ξεκινά από ``ισχυρό'' σημείο στο top layer.}
\end{quote}

% ------------------------------------------------------------
% Pre-existing (assumed) neighborhood state for layer0, layer1
% ------------------------------------------------------------
\subsubsection{Ενδεικτική κατάσταση γειτνίασης πριν το πλήρες demo (για να υπάρχουν όλες οι περιπτώσεις)}
\label{subsec:assumed_state}

Για να παρουσιαστεί pruning, θεωρούμε ότι πριν την εισαγωγή του κόμβου 5 υπάρχουν ήδη αρκετοί κόμβοι και ακμές στο layer 0.
(Το toy παράδειγμα δεν στοχεύει να αποδείξει μοναδικό γράφο, αλλά να δείξει \textbf{τη ροή των λειτουργιών}).

\paragraph{Ενδεικτικές λίστες γειτόνων στο layer 0.}
Υποθέτουμε ότι στο layer 0 ισχύουν:
\[
\begin{aligned}
N_0(3)&=\{0,4,10,7\},\\
N_0(0)&=\{1,2,3,11\},\\
N_0(4)&=\{3,6,9,8\},\\
N_0(10)&=\{3,7,9,0\},\\
N_0(7)&=\{10,9,8,3\}.
\end{aligned}
\]
Όλες οι λίστες είναι μεγέθους $\le M_0=4$.

\paragraph{Ενδεικτικές λίστες γειτόνων στο layer 1.}
Υποθέτουμε:
\[
N_1(7)=\{3,0\},\quad
N_1(3)=\{7,0\}.
\]
(Στο layer 1 ισχύει όριο $M=2$.)

% ------------------------------------------------------------
% Case 2: Full insertion demo for node 5 (level=1)
% ------------------------------------------------------------
\subsubsection{Case 2: Πλήρες demo εισαγωγής κόμβου 5 (level=1) με \texttt{search\_layer} και pruning}
\label{subsec:case_full_insert_5}

Εισάγεται ο κόμβος:
\[
5=(2.6,0.2),\quad level(5)=1.
\]
Θα παρουσιαστούν:
\begin{enumerate}
    \item Greedy descent στο top layer,
    \item \texttt{search\_layer} στο layer 1,
    \item \texttt{search\_layer} στο layer 0 με $efConstruction$,
    \item Neighbor selection,
    \item Αμφίδρομες ακμές και pruning (bounded degree).
\end{enumerate}

\paragraph{(1) Greedy descent στο layer 2.}
Το layer 2 περιέχει τους $\{0,7\}$. Ξεκινάμε από entry=0 και συγκρίνουμε:
\[
d(5,0)=2.6^2+0.2^2=6.76+0.04=6.80,
\]
\[
d(5,7)=(2.6-5)^2+(0.2-0.2)^2=2.4^2+0=5.76.
\]
Εφόσον $5.76 < 6.80$, το greedy μετακινείται στον 7, άρα entry για layer 1 είναι:
\[
ep_{L=1}=7.
\]

\paragraph{(2) \texttt{search\_layer} στο layer 1 με $efConstruction=4$.}
Ξεκινάμε από $ep=7$ και θεωρούμε $N_1(7)=\{3,0\}$.
Υπολογισμοί:
\[
d(5,3)=(2.6-3)^2+(0.2-0)^2=0.16+0.04=0.20,
\]
\[
d(5,0)=6.80.
\]
Άρα το καλύτερο σημείο στο layer 1 είναι ο 3. Με $ef=4$, ένα τυπικό σύνολο υποψηφίων είναι:
\[
W_1 \approx \{3,7,0\}.
\]

\paragraph{Neighbor selection στο layer 1 (M=2).}
Κρατάμε τους 2 καλύτερους:
\[
N_1(5)=\{3,7\},
\]
και προσθέτουμε αμφίδρομες ακμές στο layer 1:
\[
5\leftrightarrow 3,\quad 5\leftrightarrow 7 \quad (layer~1).
\]
Αν κάποιος κόμβος υπερβεί $M=2$, γίνεται pruning στο layer 1 (δεν συμβαίνει εδώ με το συγκεκριμένο toy state).

\paragraph{(3) \texttt{search\_layer} στο layer 0 με $efConstruction=4$.}
Entry για base layer λαμβάνεται το καλύτερο από το layer 1, δηλαδή:
\[
ep_{L=0}=3.
\]
Αρχικοποιούμε:
\[
C=\{(d(5,3)=0.20,3)\},\quad W=\{3\}.
\]

\textbf{Επέκταση του 3:} $N_0(3)=\{0,4,10,7\}$.
Υπολογίζουμε:
\[
d(5,4)=(2.6-3)^2+(0.2-1)^2=0.16+0.64=0.80,
\]
\[
d(5,10)=(2.6-4)^2+(0.2+0.1)^2=1.96+0.09=2.05,
\]
\[
d(5,7)=5.76,\quad d(5,0)=6.80.
\]
Με $ef=4$ κρατάμε τα 4 καλύτερα, άρα:
\[
W=\{3(0.20),4(0.80),10(2.05),7(5.76)\}.
\]

\textbf{Επέκταση του 4:} $N_0(4)=\{3,6,9,8\}$.
Υπολογίζουμε:
\[
d(5,6)=(2.6-2.7)^2+(0.2-0.9)^2=0.01+0.49=0.50,
\]
\[
d(5,9)=(2.6-4.2)^2+(0.2-0.7)^2=2.56+0.25=2.81,
\]
\[
d(5,8)=(2.6-5)^2+(0.2-1.2)^2=5.76+1=6.76.
\]
Update top-4:
\[
W=\{3(0.20),6(0.50),4(0.80),10(2.05)\}.
\]
Οι επεκτάσεις των 6 και 10 δεν προσθέτουν καλύτερους υποψηφίους (στο συγκεκριμένο toy state),
άρα τελικό:
\[
W_0=[3,6,4,10].
\]

\paragraph{(4) Neighbor selection στο layer 0 (M0=4) και diversification.}
Με $M_0=4$, το πλήθος χωράει:
\[
N_0(5)=\{3,6,4,10\}.
\]
Στο report τονίζεται ότι η heuristic επιλογή γειτόνων στο HNSW στοχεύει να αποφεύγει πλεονάζοντες γείτονες
και να κρατά ακμές που ανοίγουν διαφορετικά ``μονοπάτια''.

\paragraph{(5) Αμφίδρομες ακμές και pruning (bounded degree).}
Προστίθενται ακμές:
\[
5\leftrightarrow 3,\quad 5\leftrightarrow 6,\quad 5\leftrightarrow 4,\quad 5\leftrightarrow 10 \quad (layer~0).
\]
Τώρα ο κόμβος 3 αποκτά νέο γείτονα (τον 5). Επειδή πριν είχαμε:
\[
N_0(3)=\{0,4,10,7\},
\]
μετά γίνεται:
\[
N_0(3)=\{0,4,10,7,5\}
\]
που έχει μέγεθος 5 και υπερβαίνει $M_0=4$ $\Rightarrow$ απαιτείται pruning.

\textbf{Pruning στον κόμβο 3 (κρατάμε 4 καλύτερους).}
Υπολογίζουμε αποστάσεις από το 3:
\[
d(3,5)=0.20,\quad
d(3,4)=1,\quad
d(3,10)=(3-4)^2+(0+0.1)^2=1+0.01=1.01,
\]
\[
d(3,7)=4.04,\quad d(3,0)=9.
\]
Κρατάμε τους 4 καλύτερους:
\[
N_0(3)\leftarrow \{5,4,10,7\},
\]
και απορρίπτεται ο 0.
Αυτό δείχνει καθαρά τη λειτουργία \textbf{bounded degree}.

% ------------------------------------------------------------
% Case 3: Full query demo
% ------------------------------------------------------------
\subsubsection{Case 3: Query demo (greedy descent + base layer best-first)}
\label{subsec:case_query_demo}

Θέτουμε query:
\[
q=(2.8,0.85),\quad k=2,\quad efSearch=6.
\]

\paragraph{Greedy descent στο layer 2.}
\[
d(q,0)=2.8^2+0.85^2=7.84+0.7225=8.5625,
\]
\[
d(q,7)=(2.8-5)^2+(0.85-0.2)^2=4.84+0.4225=5.2625.
\]
Άρα entry για layer 1 είναι ο 7.

\paragraph{Greedy στο layer 1.}
Με ενδεικτικούς neighbors του 7 στο layer 1: $\{5,3\}$, υπολογίζουμε:
\[
d(q,5)=(2.8-2.6)^2+(0.85-0.2)^2=0.04+0.4225=0.4625,
\]
\[
d(q,3)=(2.8-3)^2+0.85^2=0.04+0.7225=0.7625.
\]
Ο greedy οδηγεί στο 5, άρα entry στο layer 0:
\[
ep_{L=0}=5.
\]

\paragraph{Base layer search με $efSearch=6$.}
Η best-first αναζήτηση στο base layer διατηρεί δύο δομές:
\begin{itemize}
  \item $C$ (candidates, min-heap): κόμβοι που αναμένουν επέκταση, ταξινομημένοι κατά αύξουσα απόσταση.
  \item $W$ (result set, max-heap μεγέθους $efSearch=6$): οι $ef$ καλύτεροι κόμβοι που έχουν βρεθεί ως τώρα.
\end{itemize}

\textbf{Βήμα 0 — αρχικοποίηση από $ep=5$:}
\[
d(q,5)=0.4625,\quad C=\{(0.4625,5)\},\quad W=\{5\}.
\]

\textbf{Βήμα 1 — επέκταση του 5:} $N_0(5)=\{3,6,4,10\}$.
\[
d(q,6)=0.0125,\quad d(q,4)=0.0625,\quad d(q,3)=0.7625,\quad d(q,10)=2.3425.
\]
Όλοι οι κόμβοι εισάγονται στο $W$ (είναι ακόμα $|W|<6$):
\[
W=\{6(0.0125),4(0.0625),5(0.4625),3(0.7625),10(2.3425)\}.
\]
Η συνθήκη διακοπής εξετάζει αν ο κόμβος που πρόκειται να επεκταθεί είναι μακρύτερος από τον χειρότερο στο $W$. Εδώ συνεχίζουμε.

\textbf{Βήμα 2 — επέκταση του 6:} $N_0(6) \approx \{5,4\}$ (ήδη επισκεφθέντες, παραλείπονται).

\textbf{Βήμα 3 — επέκταση του 4:} $N_0(4)=\{3,6,9,8\}$.
\[
d(q,9)=(2.8-4.2)^2+(0.85-0.7)^2=1.96+0.0225=1.9825.
\]
Ο 9 εισάγεται στο $W$:
\[
W=\{6(0.0125),4(0.0625),5(0.4625),3(0.7625),9(1.9825),10(2.3425)\}.
\]
Τώρα $|W|=6=efSearch$. Κάθε επόμενος κόμβος γίνεται δεκτός μόνο αν η απόστασή του είναι μικρότερη από τον χειρότερο ($d_{\max}=2.3425$).

\textbf{Βήμα 4 — επέκταση του 3:} $N_0(3)=\{0,4,10,7\}$. Κανένας δεν βελτιώνει το $W$ ($d(q,0)=8.56$, $d(q,7)=5.26$).

Η αναζήτηση τερματίζει όταν ο καλύτερος κόμβος στο $C$ είναι χειρότερος από τον χειρότερο στο $W$ (ή το $C$ αδειάσει). Τελικά:
\[
W=\{6(0.0125),4(0.0625),5(0.4625),3(0.7625),9(1.9825),10(2.3425)\},
\]
και τα 2 καλύτερα αποτελέσματα ($k=2$):
\[
\text{top-}k=\{6,4\}.
\]

\paragraph{Συμπέρασμα από το toy παράδειγμα.}
Το toy παράδειγμα παρουσιάζει:
\begin{itemize}
    \item \textbf{hierarchy/greedy descent} στα ανώτερα layers (γρήγορη μεταφορά κοντά στη σωστή περιοχή),
    \item \textbf{best-first \texttt{search\_layer}} στο base layer με $ef$,
    \item \textbf{neighbor selection} με bounded degree,
    \item \textbf{pruning} όταν παραβιάζεται το $M_0$,
    \item \textbf{query-time} ανάκτηση top-$k$.
\end{itemize}


\subsection{Πειραματικό Setup}
Τα πειράματα εκτελέστηκαν σε δύο σύνολα δεδομένων:
\begin{itemize}
  \item \textbf{SIFTsmall}: 10.000 διανύσματα βάσης, 100 queries, διάσταση 128 (L2 metric). Χρησιμοποιείται για γρήγορη επαλήθευση ορθότητας.
  \item \textbf{SIFT1M}: 1.000.000 διανύσματα βάσης, 10.000 queries, διάσταση 128. Χρησιμοποιείται για αξιολόγηση κλίμακας.
\end{itemize}
Και στις δύο περιπτώσεις χρησιμοποιούνται τα επίσημα ground-truth αρχεία (.ivecs) για υπολογισμό Recall@k. Όλα τα πειράματα εκτελέστηκαν σε MacBook Pro με επεξεργαστή Apple M1, 16GB RAM. Κάθε μέτρηση αποτελεί μέσο όρο 5 εκτελέσεων (με 2 warmup runs).

Συγκρίνονται τρεις υλοποιήσεις:
\begin{itemize}
  \item \textbf{HNSW NEW}: η υλοποίηση από scratch (Python/C++ wrapper).
  \item \textbf{hnswlib}: βιβλιοθήκη αναφοράς υψηλής απόδοσης.
  \item \textbf{faiss-hnsw}: υλοποίηση HNSW μέσω FAISS.
\end{itemize}

\subsection{Αποτελέσματα SIFTsmall (10K διανύσματα)}

Τα παρακάτω σχήματα (Σχήμα~\ref{fig:hnsw_comparison_plots}) παρουσιάζουν τη σύγκριση των τριών υλοποιήσεων στο SIFTsmall. Παρατηρούμε ότι η HNSW NEW επιτυγχάνει recall παρόμοιο με τις βιβλιοθήκες αναφοράς (hnswlib, faiss-hnsw), επιβεβαιώνοντας την ορθότητα της υλοποίησης. Το QPS της HNSW NEW είναι χαμηλότερο λόγω της Python overhead, ωστόσο στην C++ έκδοση τα αποτελέσματα συγκλίνουν.

\begin{figure}[H]
  \centering

  \begin{subfigure}{0.49\textwidth}
    \centering
    \includegraphics[width=\linewidth]{figures/recall_M16_efC128_k10.png}
    \caption{Recall@20 vs efSearch (M=16, efC=128)}
    \label{fig:recall_k20_m16_efc128}
  \end{subfigure}
  \hfill
  \begin{subfigure}{0.49\textwidth}
    \centering
    \includegraphics[width=\linewidth]{figures/recall_M16_efC64_k20.png}
    \caption{Recall@20 vs efSearch (M=16, efC=64)}
    \label{fig:recall_k20_m16_efc64}
  \end{subfigure}

  \vspace{0.6em}

  \begin{subfigure}{0.49\textwidth}
    \centering
    \includegraphics[width=\linewidth]{figures/qps_M8_efC128_k20.png}
    \caption{QPS vs efSearch (M=8, efC=128, k=20)}
    \label{fig:qps_k10_m16_efc64}
  \end{subfigure}
  \hfill
  \begin{subfigure}{0.49\textwidth}
    \centering
    \includegraphics[width=\linewidth]{figures/recall_vs_qps_M16_efC128_k10.png}
    \caption{Recall@10 vs QPS (M=16, efC=128)}
    \label{fig:recall_vs_qps_k10_m16_efc128}
  \end{subfigure}

  \caption{\textbf{Σχήμα 1 — SIFTsmall (10K vectors, 128-dim).} Πάνω αριστερά και πάνω δεξιά: Recall@k συναρτήσει του efSearch για δύο διαφορετικές τιμές efConstruction. Παρατηρείται ότι όλες οι υλοποιήσεις συγκλίνουν σε παρόμοιο recall, με μικρές διαφορές στα χαμηλά efSearch. Κάτω αριστερά: QPS συναρτήσει efSearch — η hnswlib παρουσιάζει υψηλότερο throughput λόγω βελτιστοποιήσεων χαμηλού επιπέδου. Κάτω δεξιά: recall-QPS tradeoff curve — μια καλή υλοποίηση πρέπει να βρίσκεται επάνω-αριστερά (υψηλό recall, υψηλό QPS).}
  \label{fig:hnsw_comparison_plots}
\end{figure}



\begin{table}[H]
\centering
\caption{Representative benchmark results for SIFTsmall (HNSW NEW vs hnswlib vs faiss-hnsw).}
\label{tab:bench_siftsmall_representative}

\begin{tabular}{lrrrrrrr}
\toprule
\textbf{Method} & \textbf{M} & \textbf{efC} & \textbf{efS} & \textbf{k} &
\textbf{Recall@k} & \textbf{QPS} & \textbf{Time (s)} \\
\midrule

\multicolumn{8}{c}{\textbf{High-quality graph regime (M=16, efC=128)}}\\
\midrule

% --- HNSW NEW ---
HNSW NEW   & 16 & 128 & 20  & 10 & 0.9760 & 20420.81 & 0.00490 \\
HNSW NEW   & 16 & 128 & 50  & 10 & 0.9970 & 10104.13 & 0.00989 \\
HNSW NEW   & 16 & 128 & 100 & 20 & 0.9990 &  6070.71 & 0.01647 \\

\midrule

% --- FAISS ---
faiss-hnsw & 16 & 128 & 20  & 10 & 0.9760 & 28042.44 & 0.00357 \\
faiss-hnsw & 16 & 128 & 50  & 10 & 0.9940 & 17396.31 & 0.00575 \\
faiss-hnsw & 16 & 128 & 100 & 10 & 1.0000 &  9480.84 & 0.01055 \\

\midrule

% --- HNSWLIB ---
hnswlib    & 16 & 128 & 20  & 10 & 0.9750 & 40808.36 & 0.00245 \\
hnswlib    & 16 & 128 & 50  & 10 & 0.9940 & 34362.95 & 0.00291 \\
hnswlib    & 16 & 128 & 100 & 20 & 0.9975 & 16569.66 & 0.00604 \\

\bottomrule
\end{tabular}
\end{table}

\subsection{Αποτελέσματα SIFT1M (1M διανύσματα)}

\textbf{Διαφορά Σχήματος 1 και Σχήματος 2:} Το Σχήμα~\ref{fig:hnsw_comparison_plots} αξιολογεί τις υλοποιήσεις στο \emph{SIFTsmall} (10.000 vectors), ενώ το Σχήμα~\ref{fig:hnsw_1m_plots} (παρακάτω) αξιολογεί στο \emph{SIFT1M} (1.000.000 vectors). Η σημαντική διαφορά είναι:
\begin{itemize}
  \item Το \textbf{SIFTsmall} δείχνει ότι η υλοποίηση είναι ορθή (recall συγκρίσιμο με αναφορά) αλλά είναι μικρό dataset — οι διαφορές QPS δεν αντικατοπτρίζουν πραγματική κλίμακα.
  \item Το \textbf{SIFT1M} αποκαλύπτει το πραγματικό κόστος κλίμακας: το QPS της HNSW NEW μειώνεται σημαντικά σε σχέση με hnswlib/faiss, αποτυπώνοντας το overhead του Python layer στο index traversal. Το recall παραμένει συγκρίσιμο, επιβεβαιώνοντας ορθότητα κατασκευής γράφου ακόμα και σε 1M vectors.
\end{itemize}

\begin{figure}[H]
  \centering

  \begin{subfigure}{0.49\textwidth}
    \centering
    \includegraphics[width=\linewidth]{1M_plots/recall_M16_efC128_k10.png}
    \caption{Recall@10 vs efSearch (M=16, efC=128)}
    \label{fig:recall_k10_m16_efc128}
  \end{subfigure}
  \hfill
  \begin{subfigure}{0.49\textwidth}
    \centering
    \includegraphics[width=\linewidth]{1M_plots/recall_M16_efC64_k20.png}
    \caption{Recall@20 vs efSearch (M=16, efC=64)}
    \label{fig:recall_k20_m16_efc64}
  \end{subfigure}

  \vspace{0.6em}

  \begin{subfigure}{0.49\textwidth}
    \centering
    \includegraphics[width=\linewidth]{1M_plots/qps_M16_efC64_k10.png}
    \caption{QPS vs efSearch (M=816 efC=64, k=10)}
    \label{fig:qps_k10_m16_efc64_1m}
  \end{subfigure}
  \hfill
  \begin{subfigure}{0.49\textwidth}
    \centering
    \includegraphics[width=\linewidth]{1M_plots/recall_vs_qps_M8_efC64_k20.png}
    \caption{Recall@10 vs QPS (M=8, efC=64)}
    \label{fig:recall_vs_qps_k20_m8_efc64_1m}
  \end{subfigure}

  \caption{\textbf{Σχήμα 2 — SIFT1M (1M vectors, 128-dim).} Σε αντίθεση με το Σχήμα 1, εδώ φαίνεται καθαρά η επίδραση της κλίμακας: η HNSW NEW εμφανίζει σημαντικά χαμηλότερο QPS (3–4$\times$) σε σχέση με hnswlib, κυρίως λόγω Python overhead ανά distance computation. Το recall παραμένει παρόμοιο και για τις τρεις υλοποιήσεις, επιβεβαιώνοντας ότι ο γράφος HNSW κατασκευάζεται ορθά. Η C++ έκδοση (HNSW NEW C++) γεφυρώνει σε μεγάλο βαθμό αυτή τη διαφορά.}
  \label{fig:hnsw_1m_plots}
\end{figure}

\begin{table}[H]
\centering
\caption{Representative benchmark results for SIFT1M (HNSW NEW vs hnswlib vs faiss-hnsw).}
\label{tab:bench_sift1m_representative}

\begin{tabular}{lrrrrrrr}
\toprule
\textbf{Method} & \textbf{M} & \textbf{efC} & \textbf{efS} & \textbf{k} &
\textbf{Recall@k} & \textbf{QPS} & \textbf{Time (s)} \\
\midrule

\multicolumn{8}{c}{\textbf{High-quality graph regime (M=16, efC=128)}}\\
\midrule

% --- HNSW NEW ---
HNSW NEW   & 16 & 128 & 20  & 10 & 0.8390 & 3243.77 & 3.0828 \\
HNSW NEW   & 16 & 128 & 50  & 10 & 0.9471 & 2377.59 & 4.2059 \\
HNSW NEW   & 16 & 128 & 100 & 10 & 0.9837 &  857.88 & 11.6567 \\

\midrule

% --- FAISS ---
faiss-hnsw & 16 & 128 & 20  & 10 & 0.8424 & 9396.84 & 1.0642 \\
faiss-hnsw & 16 & 128 & 50  & 10 & 0.9480 & 6235.89 & 1.6036 \\
faiss-hnsw & 16 & 128 & 100 & 10 & 0.9832 & 2450.09 & 4.0815 \\

\midrule

% --- HNSWLIB ---
hnswlib    & 16 & 128 & 20  & 10 & 0.8320 & 14394.57 & 0.6947 \\
hnswlib    & 16 & 128 & 50  & 10 & 0.9423 & 7242.29 & 1.3808 \\
hnswlib    & 16 & 128 & 100 & 10 & 0.9800 & 5257.88 & 1.9019 \\

\bottomrule
\end{tabular}
\end{table}

\subsection{Σύνδεση με το επόμενο milestone (DARTH ως κίνητρο)}
Το M1 καταλήγει σε καθαρή baseline υλοποίηση HNSW και αξιολόγηση.
Το επόμενο βήμα (M2) εισάγει την ιδέα του \textbf{early termination} στο base-layer search:
αντί να επιλέγεται σταθερό $efSearch$, ο αλγόριθμος θα αποφασίζει δυναμικά πότε αξίζει να συνεχίσει.

Στο πλαίσιο του DARTH, αυτό γίνεται μέσω ενός predictor που εκτιμά το τρέχον recall κατά τη διάρκεια της αναζήτησης, και τερματίζει όταν επιτευχθεί στόχος $R_t$ (target recall).

% ------------------------------------------------------------
% (Optional) Where to place figures/tables in M1
% ------------------------------------------------------------
% \subsection{Αποτελέσματα}
% \begin{table}[h]
% \centering
% \caption{Baseline σύγκριση (M1).}
% \begin{tabular}{lcccc}
% \hline
% Method & M & efC & efS & Recall@10 \\
% \hline
% HNSW-from-scratch & 16 & 100 & 100 & ... \\
% hnswlib & 16 & 100 & 100 & ... \\
% FAISS-HNSW & 16 & 100 & 100 & ... \\
% \hline
% \end{tabular}
% \end{table}



\section{Milestone M2: Ενσωμάτωση DARTH πάνω από HNSW}

\subsection{Στόχος}

Στο Milestone M1 υλοποιήθηκε ο αλγόριθμος HNSW από την αρχή και αξιολογήθηκε ως προς ορθότητα και απόδοση. 

Στο παρόν Milestone (M2) επεκτείνουμε τον HNSW με βάση την ιδέα του DARTH (Declarative Approximate Recall Target in HNSW), με στόχο:

\begin{itemize}
    \item Να επιτρέψουμε \textbf{δηλωτικό στόχο ανά query} (target recall $R_t$),
    \item Να υλοποιήσουμε \textbf{adaptive early termination} στο base layer,
    \item Να μειώσουμε το υπολογιστικό κόστος (distance computations) όταν ο στόχος recall έχει ήδη επιτευχθεί.
\end{itemize}

Σε αντίθεση με τον κλασικό HNSW όπου το $efSearch$ είναι σταθερό, εδώ το search μπορεί να τερματιστεί δυναμικά.


\subsection{Θεωρητική Αιτιολόγηση της Προσέγγισης}
Στον κλασικό HNSW, το επίπεδο ποιότητας του αποτελέσματος (recall) ελέγχεται έμμεσα μέσω της παραμέτρου $efSearch$. Μεγαλύτερη τιμή οδηγεί σε υψηλότερο recall αλλά και αυξημένο υπολογιστικό κόστος.

Ωστόσο, το $efSearch$ αποτελεί \textbf{στατική και καθολική παράμετρο}, η οποία εφαρμόζεται ομοιόμορφα σε όλα τα queries, ανεξαρτήτως της δυσκολίας τους. Στην πράξη όμως, διαφορετικά queries παρουσιάζουν διαφορετική «τοπική πολυπλοκότητα» στο γράφημα:

\begin{itemize}
\item \textbf{Εύκολα queries}: Η γειτονιά τους είναι καλά διαχωρισμένη και ο nearest neighbor εντοπίζεται γρήγορα με λίγες επεκτάσεις.
\item \textbf{Δύσκολα queries}: Βρίσκονται σε πυκνές ή αμφίσημες περιοχές του χώρου, όπου πολλές αποστάσεις είναι παρόμοιες, απαιτώντας περισσότερες επεκτάσεις για υψηλό recall.
\end{itemize}

Ως «εύκολο» query ορίζεται εκείνο για το οποίο το σωστό σύνολο $k$-NN μπορεί να προσεγγιστεί με μικρό αριθμό distance computations. Αντίθετα, «δύσκολο» query είναι εκείνο που απαιτεί μεγάλο αριθμό επεκτάσεων για να σταθεροποιηθεί το αποτέλεσμα.

Η χρήση σταθερού $efSearch$ οδηγεί σε:

\begin{itemize}
\item Υπερβολικό κόστος για εύκολα queries.
\item Ανεπαρκές recall για δύσκολα queries (αν το $efSearch$ είναι μικρό).
\end{itemize}

Η προσέγγιση DARTH εισάγει έναν \textbf{predictor recall}, ο οποίος επιτρέπει την εκτίμηση της ποιότητας κατά τη διάρκεια της αναζήτησης. Με αυτόν τον τρόπο, το stopping criterion γίνεται \textbf{δηλωτικό και προσαρμοστικό}, αντί σταθερό.

\subsection{Αρχιτεκτονική Επέκτασης}

Η βασική ιδέα είναι ότι η ιεραρχική δομή του HNSW παραμένει αμετάβλητη:

\begin{itemize}
    \item Στα ανώτερα layers εκτελείται greedy descent (beam width = 1).
    \item Η διαφοροποίηση εισάγεται αποκλειστικά στο base layer (layer 0).
\end{itemize}

Η τροποποίηση αφορά τη διαδικασία best-first search στο layer 0. Υπενθυμίζουμε ότι η best-first αναζήτηση στο layer 0 χρησιμοποιεί δύο δομές:
\begin{itemize}
  \item \textbf{Ουρά υποψηφίων $C$} (min-heap): κόμβοι που δεν έχουν επεκταθεί ακόμα, ταξινομημένοι κατά αύξουσα απόσταση από το query.
  \item \textbf{Σύνολο αποτελεσμάτων $W$} (max-heap, μέγεθος $efSearch$): οι $ef$ πλησιέστεροι κόμβοι που έχουν βρεθεί.
\end{itemize}
Σε κάθε βήμα αφαιρείται ο πλησιέστερος κόμβος από το $C$, επεκτείνεται, και οι γείτονές του εξετάζονται για εισαγωγή στο $W$. Η αναζήτηση τερματίζει όταν ο κόμβος στην κεφαλή του $C$ είναι χειρότερος από τον χειρότερο κόμβο στο $W$ (δεν μπορεί να βελτιωθεί το αποτέλεσμα). Στο DARTH, αυτή η ουρά $C$ μπορεί να τερματιστεί \emph{νωρίτερα} αν ο predictor εκτιμήσει ότι το recall έχει ήδη επιτευχθεί.

\subsection{Διαφορά από τον Κλασικό HNSW}

\begin{center}
\begin{tabular}{|l|c|c|}
\hline
 & \textbf{HNSW} & \textbf{HNSW + DARTH} \\
\hline
Termination & Όταν αδειάσει η ουρά ή φτάσει $efSearch$ & Όταν προβλεφθεί $R_p \ge R_t$ \\
Recall target & Έμμεσο (μέσω efSearch) & Δηλωτικό ($R_t$) \\
Cost control & Σταθερό efSearch & Adaptive early stop \\
\hline
\end{tabular}
\end{center}

\subsection{Feature Extraction Κατά τη Διάρκεια Αναζήτησης}

Κατά τη διάρκεια της αναζήτησης στο base layer συλλέγονται δυναμικά χαρακτηριστικά που αποτυπώνουν την κατάσταση της αναζήτησης. Αυτά χωρίζονται σε δύο ομάδες:

\textbf{Ομάδα Α — Μετρητές προόδου αναζήτησης (3 χαρακτηριστικά):}
\begin{itemize}
    \item $n_{\text{step}}$: αριθμός επεκτάσεων κόμβων (πόσοι κόμβοι αφαιρέθηκαν από το $C$)
    \item $n_{\text{dis}}$: συνολικός αριθμός distance computations
    \item $n_{\text{insert}}$: αριθμός ενημερώσεων του result set $W$ (πόσες φορές βελτιώθηκε το top-k)
\end{itemize}

\textbf{Ομάδα Β — Χαρακτηριστικά αποστάσεων (3 + 5 στατιστικά = 8 χαρακτηριστικά):}
\begin{itemize}
    \item $d_{\text{first}}$: απόσταση του αρχικού entry point από το query (πρώτη απόσταση που υπολογίστηκε)
    \item $d_{\text{closest}}$: τρέχουσα ελάχιστη απόσταση στο $W$
    \item $d_{\text{furthest}}$: απόσταση του $k$-οστού γείτονα (worst element στο $W$, δηλαδή το «πόρτα» εισόδου)
    \item Στατιστικά των αποστάσεων \emph{όλων} των κόμβων στο τρέχον $W$: mean, variance, p25, p50 (median), p75
\end{itemize}

Συνολικά: $3 + 3 + 5 = \mathbf{11}$ χαρακτηριστικά. Τα στατιστικά (mean, variance, p25, p50, p75) παρέχουν πληροφορία για τη \emph{μορφή} της κατανομής αποστάσεων στο τρέχον result set, η οποία σχετίζεται με το πόσο «σφιχτή» είναι η γειτονιά — ένας χρήσιμος δείκτης δυσκολίας του query.

\subsection{Ρόλος και Κόστος του Predictor}
Ο predictor αποτελεί εξωτερικό μοντέλο πρόβλεψης που δέχεται ως είσοδο το διάνυσμα χαρακτηριστικών και επιστρέφει την εκτιμώμενη τιμή $R_p$.

Η χρήση του συνεπάγεται:

\begin{itemize}
\item Κόστος εξαγωγής χαρακτηριστικών.
\item Κόστος inference του μοντέλου.
\end{itemize}

Επομένως, η πρόβλεψη δεν εκτελείται σε κάθε distance computation, αλλά περιοδικά. Η επιλογή του prediction interval αποτελεί κρίσιμη σχεδιαστική απόφαση, καθώς:

\begin{itemize}
\item Πολύ συχνές προβλέψεις αυξάνουν το overhead.
\item Πολύ αραιές προβλέψεις καθυστερούν τον πιθανό early termination.
\end{itemize}

Η δυναμική ρύθμιση του interval επιχειρεί να εξισορροπήσει το trade-off μεταξύ ακρίβειας πρόβλεψης και υπολογιστικού κόστους.


\subsection{Adaptive Prediction Interval}

Ο predictor έχει κόστος εκτέλεσης, άρα δεν καλείται σε κάθε distance computation. Αντ' αυτού, καλείται κάθε $p_i$ distance computations (\emph{prediction interval}). Η επιλογή του $p_i$ αποτελεί trade-off:
\begin{itemize}
  \item Μικρό $p_i$: συχνότερος έλεγχος → νωρίτερος πιθανός τερματισμός, αλλά αυξημένο overhead.
  \item Μεγάλο $p_i$: μειωμένο overhead, αλλά αργότερη αντίδραση στη σύγκλιση.
\end{itemize}

Το DARTH εισάγει \emph{adaptive} interval: αρχικά το interval είναι μεγάλο (αποφεύγουμε πρόωρη κλήση predictor όταν το result set δεν είναι ακόμα σταθερό), και μειώνεται καθώς πλησιάζουμε τον στόχο. Συγκεκριμένα, ορίζονται δύο παράμετροι:
\begin{itemize}
  \item $i_{pi}$ (\emph{initial prediction interval}): το αρχικό διάστημα — χρησιμοποιείται στην αρχή της αναζήτησης.
  \item $m_{pi}$ (\emph{minimum prediction interval}): το ελάχιστο επιτρεπτό διάστημα — δεν ελέγχουμε πυκνότερα από αυτό.
\end{itemize}

Το interval ενημερώνεται δυναμικά μετά από κάθε πρόβλεψη ως:
\[
p_i = m_{pi} + (i_{pi} - m_{pi}) \cdot (R_t - R_p)
\]
Όταν $R_p \ll R_t$: $p_i \approx i_{pi}$ (αραιοί έλεγχοι). Όταν $R_p \to R_t$: $p_i \to m_{pi}$ (πυκνοί έλεγχοι). Έτσι το σύστημα «επιταχύνει» τους ελέγχους όσο πλησιάζει στον στόχο.

\subsection{Αλγόριθμος DARTH στο Base Layer}

\textbf{Βήμα 1:} Εκκίνηση από entry point του layer 0

\textbf{Βήμα 2:} Εκτέλεση best-first traversal

\textbf{Βήμα 3:} Κάθε $p_i$ distance computations:

\begin{itemize}
    \item Υπολογισμός feature vector
    \item Πρόβλεψη $R_p$
    \item Αν $R_p \ge R_t$: τερματισμός
    \item Αλλιώς: ενημέρωση $p_i$
\end{itemize}

\textbf{Βήμα 4:} Αν δεν ικανοποιηθεί το κριτήριο, συνεχίζεται μέχρι φυσικό termination του HNSW.

\subsection{Σχεδιαστικές Αποφάσεις}
Κατά την ενσωμάτωση του DARTH λήφθηκαν οι εξής σχεδιαστικές αποφάσεις:

\begin{itemize}
\item Ο predictor εφαρμόζεται αποκλειστικά στο base layer.
\item Η ιεραρχική δομή και η διαδικασία insertion του HNSW παραμένουν αμετάβλητες.
\item Ο μηχανισμός early termination δεν επηρεάζει τη δομή του γράφου.
\item Η επέκταση είναι modular και μπορεί να απενεργοποιηθεί χωρίς αλλαγή του βασικού αλγορίθμου.
\end{itemize}

Με αυτόν τον τρόπο, η τροποποίηση αφορά αποκλειστικά το stopping criterion της best-first αναζήτησης.

\subsection{Αποτελέσματα με Dummy Predictor (Ablation Study)}
Για την αποτίμηση της επίδρασης του μηχανισμού πρόβλεψης, υλοποιήθηκε ένας dummy predictor ο οποίος επιστρέφει σταθερά $R_p = 0$.

Στην περίπτωση αυτή:

\begin{itemize}
\item Το early termination δεν ενεργοποιείται ποτέ.
\item Ο αλγόριθμος εκτελεί πλήρη αναζήτηση όπως ο κλασικός HNSW.
\end{itemize}

Τα αποτελέσματα αυτά λειτουργούν ως \textbf{baseline} για την εκτίμηση του επιπλέον κόστους που εισάγει το DARTH framework, ανεξάρτητα από την ποιότητα του predictor.

\begin{figure}[H]
  \centering

  \begin{subfigure}{0.49\textwidth}
    \centering
    \includegraphics[width=\linewidth]{DarthPaper/qps_M8_efC128_k10.png}
    \caption{Recall@10 vs efSearch (M=8, efC=128)}
    \label{fig:qps_k10_m8_efc128}
  \end{subfigure}
  \hfill
  \begin{subfigure}{0.49\textwidth}
    \centering
    \includegraphics[width=\linewidth]{DarthPaper/recall_M16_efC128_k20.png}
    \caption{Recall@20 vs efSearch (M=16, efC=128)}
    \label{fig:recall_k20_m16_efc128}
  \end{subfigure}

  \vspace{0.6em}

  \begin{subfigure}{0.49\textwidth}
    \centering
    \includegraphics[width=\linewidth]{DarthPaper/recall_M16_efC64_k20.png}
    \caption{QPS vs efSearch (M=816 efC=64, k=20)}
    \label{fig:qps_k20_m16_efc64}
  \end{subfigure}
  \hfill
  \begin{subfigure}{0.49\textwidth}
    \centering
    \includegraphics[width=\linewidth]{DarthPaper/recall_vs_qps_M8_efC128_k10.png}
    \caption{Recall@20 vs QPS (M=8, efC=64)}
    \label{fig:recall_vs_qps_k20_m8_efc128}
  \end{subfigure}

  \caption{\textbf{Σχήμα 3 — DARTH με Dummy Predictor (result set μεγέθους $k$).} Ο dummy predictor επιστρέφει πάντα $R_p=0$, άρα το early termination δεν ενεργοποιείται ποτέ — η αναζήτηση τερματίζεται φυσικά. Αυτό αποτυπώνει το \emph{καθαρό κόστος} του DARTH framework (feature extraction + predictor overhead) χωρίς όφελος από πρόωρο τερματισμό. Παρατηρείται ότι το recall του DARTH (dummy) είναι χαμηλότερο από HNSW NEW με ίδιο efSearch: αυτό οφείλεται στην επιλογή result set μεγέθους $k$ (αντί $efSearch$), η οποία κάνει την αναζήτηση πιο επιθετική αλλά λιγότερο πλήρη.}
  \label{fig:darth_dummy_plots}
\end{figure}


\begin{table}[h]
\centering
\caption{SIFTsmall – Σύγκριση DARTH (Dummy Predictor) με HNSW και Βιβλιοθήκες}
\begin{tabular}{|l|c|c|c|c|c|}
\hline
Method & M & efSearch & k & Recall@k & QPS \\
\hline
DARTH (dummy) & 8  & 20  & 10 & 0.846  & 40,048 \\
HNSW NEW      & 8  & 20  & 10 & 0.955  & 16,209 \\
faiss-hnsw    & 8  & 20  & 10 & 0.945  & 32,892 \\
hnswlib       & 8  & 20  & 10 & 0.940  & 76,287 \\
\hline
DARTH (dummy) & 16 & 50  & 20 & 0.942  & 9,116 \\
HNSW NEW      & 16 & 50  & 20 & 0.996  & 9,131 \\
faiss-hnsw    & 16 & 50  & 20 & 0.9925 & 12,550 \\
hnswlib       & 16 & 50  & 20 & 0.981  & 36,784 \\
\hline
\end{tabular}
\end{table}

\subsection{Πιθανές Πηγές Σφάλματος}
Η χρήση predictor εισάγει πιθανότητα σφάλματος:

\begin{itemize}
\item Αν το $R_p$ υπερεκτιμηθεί, η αναζήτηση μπορεί να τερματιστεί πρόωρα, μειώνοντας το πραγματικό Recall@k.
\item Αν το $R_p$ υποεκτιμηθεί, η αναζήτηση συνεχίζεται περισσότερο από όσο απαιτείται.
\end{itemize}


\subsection{Παρατήρηση Σχεδιαστικής Διαφοράς: Result Set μεγέθους $k$ έναντι $efSearch$}

Κατά την υλοποίηση του DARTH παρατηρήθηκε ότι ο ψευδοκώδικας του Algorithm~1 ορίζει το \textit{resultSet} ως heap μεγέθους $k$. 
Αυτό διαφέρει από την κλασική υλοποίηση του HNSW σε ευρέως χρησιμοποιούμενες βιβλιοθήκες (π.χ. FAISS/hnswlib), όπου στο base layer διατηρείται σύνολο υποψήφιων αποτελεσμάτων μεγέθους $efSearch$ και στο τέλος επιστρέφονται τα top-$k$.

Η επιλογή μεγέθους του result set επηρεάζει το exploration της αναζήτησης:

\begin{itemize}
    \item \textbf{Result set μεγέθους $k$ (paper-faithful)}: το κριτήριο \texttt{GetMaxDistance(resultSet)} γίνεται πιο αυστηρό νωρίς, οδηγώντας σε πιο επιθετικό pruning. Αυτό τείνει να αυξάνει το QPS αλλά μπορεί να μειώσει το Recall@k.
    \item \textbf{Result set μεγέθους $efSearch$ (standard HNSW)}: διατηρείται μεγαλύτερο σύνολο υποψηφίων κατά τη διάρκεια της αναζήτησης, ευνοώντας υψηλότερο recall με αυξημένο κόστος.
\end{itemize}

\subsection{Αποτελέσματα με result set μεγέθους $efSearch$}
\begin{figure}[H]
  \centering

  \begin{subfigure}{0.49\textwidth}
    \centering
    \includegraphics[width=\linewidth]{DarthHNSWBased/qps_M16_efC128_k10.png}
    \caption{Recall@10 vs efSearch (M=16, efC=128)}
    \label{fig:qps_k10_m16_efc128}
  \end{subfigure}
  \hfill
  \begin{subfigure}{0.49\textwidth}
    \centering
    \includegraphics[width=\linewidth]{DarthHNSWBased/recall_M8_efC128_k20.png}
    \caption{Recall@20 vs efSearch (M=8, efC=128)}
    \label{fig:recall_k20_m16_efc128}
  \end{subfigure}

  \vspace{0.6em}

  \begin{subfigure}{0.49\textwidth}
    \centering
    \includegraphics[width=\linewidth]{DarthHNSWBased/recall_M16_efC64_k20.png}
    \caption{QPS vs efSearch (M=16 efC=64, k=20)}
    \label{fig:qps_k20_m16_efc64}
  \end{subfigure}
  \hfill
  \begin{subfigure}{0.49\textwidth}
    \centering
    \includegraphics[width=\linewidth]{DarthHNSWBased/recall_vs_qps_M8_efC128_k20.png}
    \caption{Recall@20 vs QPS (M=8, efC=128)}
    \label{fig:recall_vs_qps_k20_m8_efc128}
  \end{subfigure}

  \caption{\textbf{Σχήμα 4 — DARTH με result set μεγέθους $efSearch$.} Σε αντίθεση με το Σχήμα 3, εδώ το result set κρατά $efSearch$ υποψηφίους (κλασικό HNSW). Αυτό οδηγεί σε υψηλότερο recall (ίδιο με HNSW NEW), αποδεικνύοντας ότι η διαφορά recall του Σχήματος 3 προέρχεται αποκλειστικά από την επιλογή μεγέθους result set και όχι από το DARTH mechanism. Η σύγκριση QPS δείχνει ότι ο DARTH με dummy predictor εισάγει μικρό overhead (περίπου 5--10\%) σε σχέση με plain HNSW.}
  \label{fig:darth_hnsw_plots}
\end{figure}

\begin{table}[h]
\centering
\caption{SIFTsmall – M=8, efC=64, efSearch=20, k=10}
\begin{tabular}{|l|c|c|c|c|c|}
\hline
Method & M & efSearch & k & Recall@k & QPS \\
\hline
DARTH (dummy) & 8  & 20 & 10 & 0.923  & 25,370 \\
HNSW NEW      & 8  & 20 & 10 & 0.955  & 24,872 \\
faiss-hnsw    & 8  & 20 & 10 & 0.945  & 30,730 \\
hnswlib       & 8  & 20 & 10 & 0.940  & 79,969 \\
\hline
\end{tabular}
\end{table}

\begin{table}[H]
\centering
\caption{SIFTsmall – M=8, efC=64, efSearch=50, k=10}
\begin{tabular}{|l|c|c|c|c|c|}
\hline
Method & M & efSearch & k & Recall@k & QPS \\
\hline
DARTH (dummy) & 8  & 50 & 10 & 0.955  & 12,910 \\
HNSW NEW      & 8  & 50 & 10 & 0.992  & 14,005 \\
faiss-hnsw    & 8  & 50 & 10 & 0.991  & 15,854 \\
hnswlib       & 8  & 50 & 10 & 0.991  & 49,982 \\
\hline
\end{tabular}
\end{table}

\begin{table}[H]
\centering
\caption{SIFTsmall – M=8, efC=64, efSearch=100, k=10}
\begin{tabular}{|l|c|c|c|c|c|}
\hline
Method & M & efSearch & k & Recall@k & QPS \\
\hline
DARTH (dummy) & 8  & 100 & 10 & 0.958  & 7,782 \\
HNSW NEW      & 8  & 100 & 10 & 0.998  & 8,166 \\
faiss-hnsw    & 8  & 100 & 10 & 1.000  & 9,358 \\
hnswlib       & 8  & 100 & 10 & 0.997  & 29,907 \\
\hline
\end{tabular}
\end{table}

Στο παρόν report παρουσιάζονται αποτελέσματα και για τις δύο παραλλαγές, ώστε να τεκμηριωθεί η επίδραση της σχεδιαστικής επιλογής και να εξασφαλιστεί δίκαιη σύγκριση με υλοποιήσεις αναφοράς.

\section{Milestone M3: Ενσωμάτωση PiP πάνω από HNSW}

\subsection{Στόχος}
Στο παρόν milestone μελετάται και υλοποιείται η στρατηγική \textbf{PiP (Patience in Proximity)} ως ένας απλός, predictor-free μηχανισμός early termination για την αναζήτηση στο HNSW. 
Σε αντίθεση με το DARTH, το οποίο βασίζεται σε predictor για την εκτίμηση recall, το PiP στηρίζεται αποκλειστικά στη \textbf{σταθερότητα του top-$k$ συνόλου} κατά τη διάρκεια της αναζήτησης. Στόχος είναι η μείωση του υπολογιστικού κόστους χωρίς την εισαγωγή εξωτερικού μοντέλου πρόβλεψης.

\subsection{Βασική Ιδέα}
Η βασική παρατήρηση του PiP είναι ότι, καθώς εξελίσσεται η αναζήτηση στο base layer, το σύνολο των καλύτερων γειτόνων τείνει σταδιακά να σταθεροποιείται. Αν για αρκετά συνεχόμενα βήματα το top-$k$ παραμένει σχεδόν αμετάβλητη, τότε η περαιτέρω εξερεύνηση είναι πιθανό να έχει μικρό όφελος. Έτσι, το PiP εισάγει ένα κριτήριο πρόωρου τερματισμού βασισμένο στο βαθμό επικάλυψης του top-$k$ μεταξύ διαδοχικών βημάτων.

\subsection{Μέτρο Σταθερότητας}
Έστω $N_{h,l}(q)$ το σύνολο των top-$k$ γειτόνων του query $q$ μετά από $h$ visits στο layer $l$. Τότε ο βαθμός σταθερότητας ορίζεται ως:

\[
\phi_{h,l}(q)=100 \cdot \frac{|N_{h-1,l}(q)\cap N_{h,l}(q)|}{k}
\]

Η ποσότητα $\phi_{h,l}(q)$ εκφράζει το ποσοστό των στοιχείων του top-$k$ που παραμένουν ίδια από το ένα βήμα στο επόμενο.

\subsection{Κριτήριο Early Termination}
Το PiP τερματίζει την αναζήτηση όταν η σταθερότητα του top-$k$ παραμένει υψηλή για αρκετά συνεχόμενα βήματα. Συγκεκριμένα, αν ισχυεί:
\[
\phi_{h,l}(q) \ge \gamma
\]

για $\Delta$ συνεχόμενα iterations, τότε η αναζήτηση σταματά.

Όπου:
\begin{itemize}
    \item $\gamma$: threshold σταθερότητας (π.χ. 95\%)
    \item $\Delta$: patience parameter
\end{itemize}

\subsection{Αλγόριθμος PiP στο Base Layer}

Ο αλγόριθμος PiP εκτελείται αποκλειστικά στο base layer (layer 0), ακριβώς όπως το DARTH. Τα ανώτερα layers παραμένουν αμετάβλητα (greedy descent). Η best-first αναζήτηση τροποποιείται ως εξής:

\begin{enumerate}
  \item \textbf{Εκκίνηση:} Αρχικοποίηση ουράς $C$ και result set $W$ όπως στον κλασικό HNSW. Θέτουμε μετρητή σταθερότητας $\delta \leftarrow 0$ και $W_{\text{prev}} \leftarrow \emptyset$.
  \item \textbf{Κάθε βήμα:} Επέκταση του πλησιέστερου κόμβου από $C$. Ενημέρωση $W$.
  \item \textbf{Έλεγχος σταθερότητας (μετά από κάθε επέκταση):}
    \[
    \phi = 100 \cdot \frac{|W_{\text{prev}} \cap W|}{k}
    \]
    \begin{itemize}
      \item Αν $\phi \ge \gamma$: $\delta \leftarrow \delta + 1$.
      \item Αλλιώς: $\delta \leftarrow 0$ (reset, αστάθεια).
    \end{itemize}
    Αποθήκευση $W_{\text{prev}} \leftarrow W$.
  \item \textbf{Κριτήριο τερματισμού:} Αν $\delta \ge \Delta$, τερματισμός.
  \item \textbf{Φυσικός τερματισμός:} Αν η ουρά $C$ αδειάσει, τερματισμός κανονικά.
\end{enumerate}

\subsection{Παράμετροι και Ερμηνεία}
\begin{itemize}
  \item $\gamma$ (stability threshold): ποσοστό επικάλυψης top-$k$ που θεωρείται «σταθερό». Τυπικές τιμές: 90--99\%. Χαμηλό $\gamma$ → πιο επιθετικός τερματισμός, χαμηλότερο recall.
  \item $\Delta$ (patience): αριθμός συνεχόμενων σταθερών βημάτων πριν τερματισμό. Τυπικές τιμές: 5--20. Υψηλό $\Delta$ → πιο συντηρητικός τερματισμός, υψηλότερο recall.
\end{itemize}

\subsection{Τι Θέλουμε να Πετύχουμε}
Με το PiP επιδιώκουμε:
\begin{itemize}
    \item Μείωση του αριθμού επεκτάσεων στο base layer, ιδίως για εύκολα queries όπου το top-$k$ σταθεροποιείται νωρίς.
    \item Μείωση του συνολικού υπολογιστικού κόστους χωρίς training phase ή εξωτερικό μοντέλο.
    \item Διατήρηση ικανοποιητικού Recall@k με κατάλληλη επιλογή $(\gamma, \Delta)$.
\end{itemize}

\textbf{Σύγκριση με DARTH:} Το PiP δεν απαιτεί training phase ούτε predictor, άρα είναι εύκολα εφαρμόσιμο σε νέα datasets. Ωστόσο, δεν παρέχει δηλωτικό recall target ($R_t$) ανά query — ο χρήστης ελέγχει την ποιότητα έμμεσα μέσω $(\gamma, \Delta)$.

\section{Milestone M4: Ενσωμάτωση Ada-ef πάνω από HNSW}

Στο Milestone M4 εξετάζεται η προσέγγιση \textbf{Ada-ef (Adaptive-ef)}, η οποία επιλέγει δυναμικά την τιμή του exploration factor $e$ για κάθε query, με στόχο να προσεγγίσει ένα επιθυμητό επίπεδο recall με όσο το δυνατόν μικρότερο υπολογιστικό κόστος.
Σε αντίθεση με τον κλασικό HNSW, όπου το $efSearch$ είναι σταθερό για όλα τα queries, το Ada-ef εκτίμα online τη δυσκολία κάθε query και αντιστοιχίζει κατάλληλο $ef$
\subsection{Βασική Ιδέα}
Η βασική ιδέα του Ada-ef είναι ότι διαφορετικά queries απαιτούν διαφορετικό βαθμό εξερεύνησης στο base layer. Ορισμένα queries μπορούν να πετύχουν υψηλό recall με μικρό $ef$, ενώ άλλα απαιτούν σημαντικά μεγαλύτερο exploration budget.

Το Ada-ef λειτουργεί σε δύο φάσεις:
\begin{itemize}
    \item \textbf{Offline phase}: υπολογίζονται στατιστικά του dataset και κατασκευάζεται ένας πίνακας εκτίμησης $ef$.
    \item \textbf{Online phase}: συλλέγονται αποστάσεις κατά την αναζήτηση στο base layer, υπολογίζεται score δυσκολίας για το query, και επιλέγεται δυναμικά το κατάλληλο $ef$.
\end{itemize}

\subsection{Θεωρητική Βάση}

Η προσέγγιση Ada-ef βασίζεται στην παρατήρηση ότι για high dimensional learned embeddings, η πλήρης λίστα αποστάσεων μεταξύ ενός query $q$ και του dataset $V$ μπορεί να προσεγγιστεί από κανονική κατανομή:
\[
FDL(q,V) \approx \mathcal{N}(\mu_q,\sigma_q^2)
\]
Η κατανομή αυτή εκτιμάται αποδοτικά με χρήση dataset-level στατιστικών και ενός μικρού δείγματος αποστάσεων που συλλέγεται online.


\subsection{Query Scoring}
Κατά το online στάδιο συλλέγεται μία λίστα αποστάσεων $D$ από κόμβους του base layer. Έπειτα εκτιμάται ένα ποσοτημόριο της απόστασης και ορίζεται score δυσκολίας:
\[
s(q)=\left|\left\{d \in D \;:\; d \le Q_p(\hat{FDL}(q,V))\right\}\right|
\]
όπου $Q_p(\hat{FDL}(q,V))$ είναι το εκτιμώμενο $p$-ποσοστημόριο της distance distribution.
Το score αυτό χρησιμοποιείται για να εκτιμηθεί αν το query είναι εύκολο ή δύσκολο:
\begin{itemize}
    \item Μεγάλο $s(q)$: πολλές μικρές αποστάσεις, άρα πιθανώς μικρό $ef$ αρκεί
    \item Μικρό $s(q)$: πιο δύσκολο query, άρα απαιτείται μεγαλύτερο $ef$
\end{itemize}

\subsection{Δυναμική επιλογή του \texorpdfstring($ef$){ef}}
Με βάση το score και το στόχο recall $R_t$, το Ada-ef επιλέγει δυναμικά exploration factor μέσω ενός offline πίνακα αντιστοίχισης:
\[
ef = T(s(q),R_t)
\]
όπου $Τ$ είναι ο ef-estimation table που έχει προυπολογιστεί offline.

\subsection{Αλγόριθμος Ada-ef}

\textbf{Offline φάση (εκτελείται μία φορά πριν την αναζήτηση):}
\begin{enumerate}
  \item Δειγματοληψία $n_s$ τυχαίων διανυσμάτων από το dataset.
  \item Υπολογισμός αποστάσεων μεταξύ των δειγμάτων, εκτίμηση dataset-level statistics $(\mu_D, \sigma_D)$.
  \item Για κάθε δυσκολία $s$ και στόχο recall $R_t$: εκτέλεση αναζήτησης με διαφορετικά $ef$ και καταγραφή του ελάχιστου $ef$ που επιτυγχάνει $R_t$. Αυτό κατασκευάζει τον πίνακα $T(s, R_t) \to ef$.
\end{enumerate}

\textbf{Online φάση (για κάθε query $q$):}
\begin{enumerate}
  \item Εκτέλεση σύντομης αρχικής αναζήτησης (μικρό $ef_0$) για συλλογή δείγματος αποστάσεων $D$ από το base layer.
  \item Εκτίμηση της κατανομής $\widehat{FDL}(q,V) \approx \mathcal{N}(\hat{\mu}_q, \hat{\sigma}_q^2)$ χρησιμοποιώντας τα dataset-level statistics και τις αποστάσεις $D$.
  \item Υπολογισμός score δυσκολίας:
  \[
  s(q) = \left|\left\{d \in D : d \le Q_p\bigl(\widehat{FDL}(q,V)\bigr)\right\}\right|
  \]
  Ένα υψηλό $s(q)$ σημαίνει ότι πολλές αποστάσεις είναι μικρές → εύκολο query (μικρό $ef$ αρκεί).
  \item Αναζήτηση $ef$ από τον offline πίνακα: $ef^* = T(s(q), R_t)$.
  \item Εκτέλεση πλήρους αναζήτησης HNSW με $efSearch = ef^*$.
\end{enumerate}

\subsection{Παράμετροι}
\begin{itemize}
  \item $R_t$: δηλωτικός στόχος recall (ίδιος ρόλος με το DARTH).
  \item $n_b$ (\emph{bins}): αριθμός bins δυσκολίας στον offline πίνακα. Τυπικά 100.
  \item $\delta$ (\emph{recall tolerance}): επιτρεπτή απόκλιση recall κατά τη δημιουργία του πίνακα. Τυπικά 0.05.
  \item $p$ (\emph{percentile}): το ποσοστημόριο $p$ που χρησιμοποιείται στον υπολογισμό $Q_p$.
\end{itemize}

\subsection{Τι θέλουμε να πετύχουμε}
Με το Ada-ef επιδιώκουμε:

\begin{itemize}
    \item \textbf{Προσαρμοστικό $ef$ ανά query}: εύκολα queries λαμβάνουν μικρό $ef$, δύσκολα μεγαλύτερο — αποφεύγοντας τόσο over-searching όσο και under-searching.
    \item \textbf{Declarative recall target}: ο χρήστης ορίζει $R_t$ και το σύστημα επιλέγει αυτόματα το κατάλληλο $ef$.
    \item \textbf{Αποφυγή training recall predictor}: αντί machine learning, χρησιμοποιείται distribution-aware rule-based μηχανισμός.
\end{itemize}

\textbf{Σύγκριση με DARTH:} Το Ada-ef δεν παρακολουθεί τη διαδικασία αναζήτησης βήμα-βήμα (δεν εφαρμόζει early termination εντός της αναζήτησης), αλλά επιλέγει $ef$ \emph{πριν} ξεκινήσει η αναζήτηση. Αυτό είναι ταχύτερο υπολογιστικά αλλά λιγότερο ευέλικτο: αν η εκτίμηση δυσκολίας είναι ανακριβής, το $ef$ μπορεί να είναι υπερ- ή υπο-αρκετό.

\textbf{Σύγκριση με PiP:} Το Ada-ef παρέχει declarative recall target (όπως DARTH), ενώ το PiP όχι. Ωστόσο, το PiP δεν απαιτεί offline φάση.
\end{document}



